% GENERATED FILE. Edit canonical lessons, then run npm run generate:books.
% canonical-source-hash: fnv1a64:091d4eb91efb0bdd
% canonical-lessons: FR-C16-suis, FR-C16-es, FR-C16-est, FR-C16-sommes, FR-C16-etes, FR-C16-sont, FR-C16-esse, FR-C16-fu, FR-C16-stare, FR-C16-suppletion, FR-C16-ser-estar, FR-C16-practice

\chapter{To Be, and the Three Verbs Inside It}
\label{ch:etre}
\hlchaptermodality{19}

\begin{chapteropening}
\textbf{By the end of this chapter:} \emph{I can use all six forms of être, and explain why it has no pattern to find.}

The most irregular verb in French, and the irregularity turns out to have a shape. There is no pattern in être because être is THREE verbs: es- from esse for the present and infinitive, fu- from fuī for the passé simple — the same root as English be — and ét- from stāre, 'to stand', which supplied a whole limb of participles and the imperfect. The patching is called suppletion, and English does it too, with go and went. Spanish inherited the same two Latin verbs and made the opposite choice: it kept ser and estar apart and made you pick.
\end{chapteropening}

\section[je suis]{je suis — I am}

\label{lesson:FR-C16-suis}



\begin{quote}
The partner of \emph{avoir}, and the most irregular verb in French. That irregularity is not carelessness; it is age.
\end{quote}

\begin{sounds}
\begin{itemize}
\raggedright
  \item \textbf{je suis} = \emph{zhuh SWEE}. The \textbf{s} at the end is silent.
  \item The \textbf{ui} is a single glide: start at the tight \textbf{u} of \emph{tu} and slide straight into \emph{ee}.
  \item Do not hear an English \emph{swiss}; there is no \emph{s} at the end.
\end{itemize}
\end{sounds}

\begin{cousinweb}
\textbf{suis} is Latin \textbf{sum}, \textquotedblleft{}I am.\textquotedblright{}

Before you go looking for a pattern in this verb: \textbf{there is not one.} \emph{être} is the most irregular verb in French, and it is irregular for the same reason English \emph{be / am / is / was} is. These are among the \textbf{oldest word-shapes in Europe}, said millions of times a day for three thousand years, and constant use does not regularise a word — it protects it. Rare verbs get tidied up by analogy; the commonest ones are said too often for anyone to forget them, so they keep whatever shape history left them with.

So the six forms are six things to learn, and that is the honest situation. What this chapter can give you instead of a pattern is an \textbf{explanation} — because the irregularity turns out to have a shape once you know where the pieces came from.
\end{cousinweb}

\subsection*{Your turn}
Say these aloud:

\begin{itemize}
\raggedright
  \item \textquotedblleft{}je suis\textquotedblright{} — \emph{zhuh SWEE}, the final \emph{s} silent
  \item a whole sentence — \textquotedblleft{}Je suis fatigué\textquotedblright{}
  \item the Latin — \textquotedblleft{}sum, suis\textquotedblright{}
\end{itemize}

\begin{tcolorbox}[breakable,colback=teal!4,colframe=teal!35!black,title={Before you move on}]
How do you say \textquotedblleft{}I am\textquotedblright{}? (\textbf{je suis}, \emph{zhuh SWEE}.) Is the final \textbf{s} pronounced? (\textbf{No.}) What Latin word is it? (\emph{\textbf{sum}}.) Why is this verb so irregular? (It is \textbf{very old and very common} — constant use protects a shape rather than tidying it.) Next: the same verb, to a friend.
\end{tcolorbox}
\section[tu es]{tu es — you are}

\label{lesson:FR-C16-es}



\begin{quote}
After \emph{suis}, the shortest form in the verb — and one that is easy to swallow entirely.
\end{quote}

\begin{sounds}
\begin{itemize}
\raggedright
  \item \textbf{tu es} = \emph{tü ay}. Two letters, one vowel sound.
  \item The \textbf{s} is silent, as it was in \textbf{tu as}.
  \item Careful: \textbf{tu es} and \textbf{tu as} differ only in that vowel — \emph{ay} against \emph{ah}. Say them back to back until the difference is deliberate.
\end{itemize}
\end{sounds}

\begin{cousinweb}
\textbf{es} is Latin \textbf{es}, unchanged. Latin said \emph{es} for \textquotedblleft{}you are\textquotedblright{} and French still does, which for a form this heavily used is unusual — most such words are ground down past recognition.

The pair to watch is not Latin, though. It is \textbf{tu es} against \textbf{tu as}: \textquotedblleft{}you are\textquotedblright{} and \textquotedblleft{}you have,\textquotedblright{} one vowel apart, both extremely common, both with a silent \textbf{s}. English keeps them far apart; French does not.

That is the practical difficulty of this verb. Not the etymology — the fact that its forms are short, unstressed, and easily confused with \emph{avoir}'s.
\end{cousinweb}

\subsection*{Your turn}
Say these aloud:

\begin{itemize}
\raggedright
  \item \textquotedblleft{}tu es\textquotedblright{} — \emph{tü ay}
  \item the pair to keep apart — \textquotedblleft{}tu es, tu as\textquotedblright{}
  \item both verbs together — \textquotedblleft{}je suis, j'ai\textquotedblright{}
\end{itemize}

\begin{tcolorbox}[breakable,colback=teal!4,colframe=teal!35!black,title={Before you move on}]
How do you say \textquotedblleft{}you are\textquotedblright{} to a friend? (\textbf{tu es}, \emph{tü ay}.) Which other very common form does it nearly collide with? (\emph{\textbf{tu as}}, \textquotedblleft{}you have\textquotedblright{}.) What is the difference? (Only the \textbf{vowel} — \emph{ay} against \emph{ah}.) What Latin word is \emph{es}? (\emph{\textbf{es}} — unchanged.) Next: he and she.
\end{tcolorbox}
\section[il est, elle est]{il est, elle est — he is, she is}

\label{lesson:FR-C16-est}



\begin{quote}
You have said this one already, many times, without knowing it was part of a verb.
\end{quote}

\begin{sounds}
\begin{itemize}
\raggedright
  \item \textbf{il est} = \emph{ee-LAY}. The \textbf{s} and \textbf{t} are both silent, and the \textbf{l} of \emph{il} runs into the vowel.
  \item Before a vowel the buried \textbf{t} wakes: \textbf{est-il ?} = \emph{eh-TEEL}.
  \item \textbf{elle est} = \emph{el-AY}.
\end{itemize}
\end{sounds}

\begin{cousinweb}
\textbf{est} is Latin \textbf{est}, also unchanged — and it is very probably the single most frequently said word-form in French.

You already own it. \textbf{il est trois heures}, \textquotedblleft{}it is three o'clock,\textquotedblright{} is this exact form; so is every time you have said what time it was. The clock chapter handed you a piece of \emph{être} several chapters before the verb arrived, which is the ordinary way a language teaches itself: the commonest forms turn up inside fixed phrases long before anybody explains them.

Note also that \textbf{il} here is not \textquotedblleft{}he.\textquotedblright{} In \emph{il est trois heures} it refers to nothing at all — French requires a subject even when there is nobody to be one. That is the same reason the pronoun cannot be dropped: the sentence needs the slot filled.
\end{cousinweb}

\subsection*{Your turn}
Say these aloud:

\begin{itemize}
\raggedright
  \item \textquotedblleft{}il est, elle est\textquotedblright{}
  \item what you already knew — \textquotedblleft{}il est trois heures\textquotedblright{}
  \item the three singulars — \textquotedblleft{}je suis, tu es, il est\textquotedblright{}
\end{itemize}

\begin{tcolorbox}[breakable,colback=teal!4,colframe=teal!35!black,title={Before you move on}]
How do you say \textquotedblleft{}he is\textquotedblright{}? (\textbf{il est}, \emph{ee-LAY}.) Where have you said this form before? (\textbf{Telling the time} — \emph{il est trois heures}.) Who is \emph{il} in that sentence? (\textbf{Nobody} — French needs the subject slot filled even when there is no subject.) Next: the plural, and a form with no French shape at all.
\end{tcolorbox}
\section[nous sommes]{nous sommes — we are}

\label{lesson:FR-C16-sommes}



\begin{quote}
Every verb you have met marks \emph{nous} with \textbf{-ons}. Watch what this one does.
\end{quote}

\begin{sounds}
\begin{itemize}
\raggedright
  \item \textbf{nous sommes} = \emph{noo SOM}. The \textbf{s} of \emph{nous} buzzes across into the vowel.
  \item The final \textbf{-es} is silent; the word ends on the \textbf{m}.
  \item Compare \textbf{nous avons}, which ends in a nasal. This one does not.
\end{itemize}
\end{sounds}

\begin{cousinweb}
\textbf{sommes} is Latin \textbf{sumus}.

And notice what is missing. Every other verb you have met marks \emph{nous} with \textbf{-ons}: \emph{nous avons}, \emph{nous parlons}. This verb does not. There is no \emph{-ons} here at all, and no way to derive \emph{sommes} from anything.

That is the shape of this verb's irregularity: it does not bend the rules, it predates them. The \textbf{-ons} ending is a regularity French built up over centuries, and \emph{être} was already too well-worn and too heavily used to be dragged into it.

So \emph{sommes} is not an exception to a rule. It is a survival from before the rule existed.
\end{cousinweb}

\subsection*{Your turn}
Say these aloud:

\begin{itemize}
\raggedright
  \item \textquotedblleft{}nous sommes\textquotedblright{} — \emph{noo SOM}
  \item the missing ending — \textquotedblleft{}nous avons, nous parlons, nous sommes\textquotedblright{}
  \item the reason — \textquotedblleft{}it predates the rule, it does not break it\textquotedblright{}
\end{itemize}

\begin{tcolorbox}[breakable,colback=teal!4,colframe=teal!35!black,title={Before you move on}]
How do you say \textquotedblleft{}we are\textquotedblright{}? (\textbf{nous sommes}, \emph{noo SOM}.) Which usual ending is missing? (\textbf{-ons}.) Why? (\emph{être} is \textbf{older than that regularity} and was too common to be pulled into it.) What Latin word is it? (\emph{\textbf{sumus}}.) Next: the form that keeps a hat.
\end{tcolorbox}
\section[vous êtes]{vous êtes — you are}

\label{lesson:FR-C16-etes}



\begin{quote}
The politeness form, and the one that wears its history on the page.
\end{quote}

\begin{sounds}
\begin{itemize}
\raggedright
  \item \textbf{vous êtes} = \emph{voo-ZET}. The \textbf{s} of \emph{vous} buzzes into the vowel.
  \item \texttt{circumflex} — the hat on \textbf{ê} marks a \textbf{lost s}, the same gravestone you met on \textbf{août} and \textbf{âge}.
  \item The final \textbf{-es} is silent.
\end{itemize}
\end{sounds}

\begin{cousinweb}
\textbf{êtes} is Latin \textbf{estis}, \textquotedblleft{}you are\textquotedblright{} — and you can read the erosion straight off the page.

\emph{estis} $\to$ \emph{estes} $\to$ \textbf{êtes}. The \textbf{s} after the first vowel fell silent and then fell out of the spelling, and the circumflex was put there to mark where it had been. This is the third time you have met that hat doing exactly this job, and by now it should be a reflex: \textbf{see a circumflex, try putting an \emph{s} back.}

Do it here and \emph{êtes} becomes \emph{estes}, which is visibly Latin \emph{estis}. Do it on \textbf{août} and you get \emph{aoust}. Do it on \textbf{âge} and you get \emph{aage}, from \emph{aetāticum}.

The infinitive works the same way: \textbf{être} was \emph{estre}, from Vulgar Latin \emph{*essere}. The hat is not decoration on any of them.
\end{cousinweb}

\subsection*{Your turn}
Say these aloud:

\begin{itemize}
\raggedright
  \item \textquotedblleft{}vous êtes\textquotedblright{} — \emph{voo-ZET}
  \item put the \emph{s} back — \textquotedblleft{}êtes, estes, estis\textquotedblright{}
  \item the same trick on the infinitive — \textquotedblleft{}être, estre\textquotedblright{}
  \item the two plurals — \textquotedblleft{}nous sommes, vous êtes\textquotedblright{}
\end{itemize}

\begin{tcolorbox}[breakable,colback=teal!4,colframe=teal!35!black,title={Before you move on}]
How do you say \textquotedblleft{}you are\textquotedblright{} politely? (\textbf{vous êtes}, \emph{voo-ZET}.) What does the circumflex mark? (A \textbf{lost s}.) Put it back — what do you get? (\emph{\textbf{estes}}, from Latin \emph{estis}.) And on the infinitive \emph{être}? (\emph{\textbf{estre}}.) Next: the last form, and the one you must not confuse with \emph{avoir}.
\end{tcolorbox}
\section[ils sont, elles sont]{ils sont, elles sont — they are}

\label{lesson:FR-C16-sont}



\begin{quote}
The sixth form, and the one you were warned about two chapters ago.
\end{quote}

\begin{sounds}
\begin{itemize}
\raggedright
  \item \textbf{ils sont} = \emph{eel SO(N)}. The \textbf{on} is nasal; the \textbf{t} is silent.
  \item Here is the pair to drill: \textbf{ils sont} \emph{hisses}, \textbf{ils ont} \emph{buzzes}. \emph{eel-\textbf{s}o(n)} against \emph{eel-\textbf{z}o(n)}.
  \item That single consonant is the whole difference between \textbf{they are} and \textbf{they have}.
\end{itemize}
\end{sounds}

\begin{cousinweb}
\textbf{sont} is Latin \textbf{sunt}.

You now have both of the verbs French is built on, and their \emph{ils} forms are one sound apart. That is worth two minutes of drilling on its own, because the two sentences it distinguishes are both things you will say constantly.

The full set, and you built every piece of it:

\textbf{je suis · tu es · il est · nous sommes · vous êtes · ils sont}

Six forms and no pattern — which is what the next three lessons are about, because the \emph{reason} there is no pattern turns out to be more interesting than a pattern would have been.
\end{cousinweb}

\subsection*{Your turn}
Say these aloud:

\begin{itemize}
\raggedright
  \item \textquotedblleft{}ils sont\textquotedblright{} — \emph{eel SO(N)}, hissing
  \item the pair — \textquotedblleft{}ils sont, ils ont\textquotedblright{}; hiss for are, buzz for have
  \item all six — \textquotedblleft{}je suis, tu es, il est, nous sommes, vous êtes, ils sont\textquotedblright{}
\end{itemize}

\begin{tcolorbox}[breakable,colback=teal!4,colframe=teal!35!black,title={Before you move on}]
How do you say \textquotedblleft{}they are\textquotedblright{}? (\textbf{ils sont}, \emph{eel SO(N)}.) How do you hear it against \emph{ils ont}? (\textbf{A hiss} against a \textbf{buzz}.) Give all six forms of \emph{être}. Is there a pattern in them? (\textbf{No} — and the next lessons explain why.) Next: where these forms actually came from.
\end{tcolorbox}
\section[esse]{esse — the first stem}

\label{lesson:FR-C16-esse}



\begin{quote}
Six forms with no pattern between them. Put them beside their Latin and something appears — not a pattern, but a \textbf{seam}.
\end{quote}

\begin{cousinweb}
Every present form of \emph{être} comes from Latin \textbf{esse}, \textquotedblleft{}to be\textquotedblright{}:

\begin{itemize}
\raggedright
  \item \textbf{je suis} $\leftarrow$ \emph{sum} · \textbf{nous sommes} $\leftarrow$ \emph{sumus} · \textbf{ils sont} $\leftarrow$ \emph{sunt}
  \item and the two short ones, \textbf{tu es} $\leftarrow$ \emph{es}, \textbf{il est} $\leftarrow$ \emph{est}, arrived practically untouched
\end{itemize}

Say each pair aloud and the irregularity stops looking random. The \textbf{s-} is there every time, and the forms are \textbf{as scattered in Latin as they are in French} — \emph{sum, es, est, sumus, estis, sunt} is not a tidy paradigm either. French did not break this verb. It \textbf{received} it broken, and Latin had received it broken too.

The infinitive comes from the same place, by a detour. Classical Latin's infinitive was \emph{esse}, which is short and odd; Vulgar Latin rebuilt it on the ordinary pattern as \textbf{*essere}, and \emph{*essere} $\to$ \emph{estre} $\to$ \textbf{être}.

So one stem, \textbf{es-}, accounts for the present tense and the infinitive. That is most of what you have learned so far. The rest of the verb comes from somewhere else entirely.
\end{cousinweb}

\subsection*{Your turn}
Say these aloud:

\begin{itemize}
\raggedright
  \item the pairs — \textquotedblleft{}sum suis, sumus sommes, sunt sont\textquotedblright{}
  \item the stem — \textquotedblleft{}es-, for the present and the infinitive\textquotedblright{}
  \item the rebuild — \textquotedblleft{}esse, essere, estre, être\textquotedblright{}
\end{itemize}

\begin{tcolorbox}[breakable,colback=teal!4,colframe=teal!35!black,title={Before you move on}]
Which Latin verb are all the present forms from? (\emph{\textbf{esse}}.) Was Latin's version regular? (\textbf{No} — \emph{sum, es, est, sumus, estis, sunt} is already scattered; French inherited the mess.) Where does the infinitive \emph{être} come from? (Vulgar Latin \textbf{*essere}, a rebuild of \emph{esse}.) Next: a part of the verb that is not from \emph{esse} at all.
\end{tcolorbox}
\section[je fus]{je fus — the second stem}

\label{lesson:FR-C16-fu}



\begin{quote}
The past of this verb does not begin with an \textbf{s}. That is not a sound change. It is a different verb.
\end{quote}

\begin{cousinweb}
The \emph{passé simple} of \emph{être} — the literary past you learned to recognise — is \textbf{je fus}, \textbf{il fut}.

There is no way to get \emph{fus} from \emph{suis}, and nobody has ever tried. It comes from Latin \textbf{fuī}, and \emph{fuī} comes from Proto-Indo-European \textbf{*b\textsuperscript{h}uH-}, \textquotedblleft{}to grow, to become.\textquotedblright{}

That root went somewhere you will recognise. \textbf{*b\textsuperscript{h}uH-} is also the source of English \textbf{be}. So English's \emph{be} and French's \emph{fus} are the same ancient word, and neither of them is related to French \emph{suis} — which is the cousin of English \emph{is} and \emph{am}.

English has exactly the same patchwork, and nobody notices because it is their own: \textbf{am} and \textbf{is} from one root, \textbf{be} from another, \textbf{was} and \textbf{were} from a third. French is doing what English does. It is only visible here because it is somebody else's language.
\end{cousinweb}

\subsection*{Your turn}
Say these aloud:

\begin{itemize}
\raggedright
  \item the form — \textquotedblleft{}je fus, il fut\textquotedblright{}
  \item the root — \textquotedblleft{}fuī, from *b\textsuperscript{h}uH-: to grow, to become\textquotedblright{}
  \item the surprise — \textquotedblleft{}that root is English \emph{be}\textquotedblright{}
  \item the English patchwork — \textquotedblleft{}am, is, be, was: three roots\textquotedblright{}
\end{itemize}

\begin{tcolorbox}[breakable,colback=teal!4,colframe=teal!35!black,title={Before you move on}]
What is \emph{je fus}? (The \textbf{passé simple} of \emph{être} — \textquotedblleft{}I was\textquotedblright{}.) Is it related to \emph{suis}? (\textbf{No} — different root entirely.) Which Latin word is it from? (\emph{\textbf{fuī}}.) Which English word comes from the same root? (\textbf{be}.) Does English do the same thing? (\textbf{Yes} — \emph{am}/\emph{is}, \emph{be}, and \emph{was} come from \textbf{three} different roots.) Next: the third stem, and it is a verb about standing.
\end{tcolorbox}
\section[été]{été — the third stem}

\label{lesson:FR-C16-stare}



\begin{quote}
Two stems so far. The third is the largest, and it is not a form of \textquotedblleft{}to be\textquotedblright{} in any language — it is a verb about \textbf{standing}.
\end{quote}

\begin{cousinweb}
The past participle of \emph{être} is \textbf{été}. It has nothing to do with \emph{suis} and nothing to do with \emph{fus}. It comes from Latin \textbf{stāre}, \textquotedblleft{}\textbf{to stand}.\textquotedblright{}

And it is not one stray form. \textbf{Every} \emph{être} form beginning \textbf{ét-} comes from the standing verb:

\begin{itemize}
\raggedright
  \item \textbf{été} — the past participle $\leftarrow$ \emph{stātum}
  \item \textbf{étant} — the present participle $\leftarrow$ \emph{stantem}
  \item \textbf{j'étais}, \textbf{il était} — the imperfect $\leftarrow$ \emph{stābam}
\end{itemize}

That is a whole limb of the verb, borrowed wholesale from a different one.

The idea is not far-fetched. \textquotedblleft{}How are you?\textquotedblright{} is \emph{how do you stand?} in more languages than not — English still says \emph{how do you stand on this}, and \emph{standing} means status. A verb for standing drifting into a verb for being is one of the best-worn paths in language.

\emph{Stāre} left French other things too, outside \emph{être}: \textbf{rester} $\leftarrow$ \emph{re-stāre}, \textquotedblleft{}to stay,\textquotedblright{} and \textbf{coûter} $\leftarrow$ \emph{constāre}, \textquotedblleft{}to cost.\textquotedblright{} What it did \textbf{not} do in French is survive as a separate verb meaning \textquotedblleft{}to be.\textquotedblright{}
\end{cousinweb}

\subsection*{Your turn}
Say these aloud:

\begin{itemize}
\raggedright
  \item the participle — \textquotedblleft{}été\textquotedblright{}
  \item the limb — \textquotedblleft{}été, étant, j'étais\textquotedblright{}: all from \emph{stāre}
  \item the meaning — \textquotedblleft{}to stand\textquotedblright{}
  \item its other children — \textquotedblleft{}rester, coûter\textquotedblright{}
\end{itemize}

\begin{tcolorbox}[breakable,colback=teal!4,colframe=teal!35!black,title={Before you move on}]
What is the past participle of \emph{être}? (\textbf{été}.) Which Latin verb is it from? (\emph{\textbf{stāre}}, \textquotedblleft{}to \textbf{stand}\textquotedblright{}.) Which other \emph{être} forms come from the same place? (Every one beginning \textbf{ét-} — \emph{étant}, \emph{j'étais}, \emph{il était}.) Name two French verbs outside \emph{être} from \emph{stāre}. (\textbf{rester}, \textbf{coûter}.) Next: the name for what this verb is doing.
\end{tcolorbox}
\section[suppletion]{suppletion — three verbs in a trench coat}

\label{lesson:FR-C16-suppletion}



\begin{quote}
You have three stems and three different origins. There is a name for this, and having the name makes it stop being strange.
\end{quote}

\begin{grammarlens}[title={Grammar Lens: a paradigm filled in from elsewhere}]
Lay out what \emph{être} is made of:

\begin{itemize}
\raggedright
  \item \textbf{es-} — the present tense and the infinitive, from \emph{esse}
  \item \textbf{fu-} — the \emph{passé simple}, from \emph{fuī}
  \item \textbf{ét-} — the participles and the imperfect, from \emph{stāre}
\end{itemize}

One verb, three ancestors, and the joins are invisible unless somebody points at them.

The name for this is \textbf{suppletion}: one verb's paradigm filled in with pieces of another. It is not a French disease and not a sign of decay. It happens to the \textbf{most used verbs in every language}, and only to them — a verb has to be extremely common for speakers to tolerate a completely unpredictable form, and extremely common verbs are also the ones people reach for when another verb is missing a form.

English does it constantly and you have never thought about it:

\begin{itemize}
\raggedright
  \item \textbf{go} / \textbf{went} — \emph{went} was stolen outright from \emph{wend}
  \item \textbf{am} / \textbf{is} / \textbf{be} / \textbf{was} — the same three-way patchwork as \emph{être}
  \item \textbf{good} / \textbf{better} — \emph{better} is not from \emph{good} at all
\end{itemize}

So the answer to \textquotedblleft{}why is \emph{être} so irregular?\textquotedblright{} is now complete. It is irregular because it is old, common, and \textbf{assembled from three verbs}. There was never a pattern to find, and the time you might have spent hunting for one is better spent on the six forms themselves.
\end{grammarlens}

\subsection*{Your turn}
Say these aloud:

\begin{itemize}
\raggedright
  \item the three stems — \textquotedblleft{}es-, fu-, ét-\textquotedblright{}
  \item the word — \textquotedblleft{}suppletion\textquotedblright{}
  \item the English case — \textquotedblleft{}go, went; and went was taken from wend\textquotedblright{}
  \item the conclusion — \textquotedblleft{}no pattern, because it is three verbs\textquotedblright{}
\end{itemize}

\begin{tcolorbox}[breakable,colback=teal!4,colframe=teal!35!black,title={Before you move on}]
Name \emph{être}'s three stems and what each covers. (\textbf{es-} present and infinitive, \textbf{fu-} the \emph{passé simple}, \textbf{ét-} participles and imperfect.) What is the name for a paradigm filled in from another verb? (\textbf{Suppletion}.) Give an English example. (\emph{\textbf{go\emph{/}went}}, or \emph{am}/\emph{be}/\emph{was}, or \emph{good}/\emph{better}.) Which verbs does it happen to? (The \textbf{most used ones} — only they are common enough to survive it.) Next: what Spanish did with the same two verbs.
\end{tcolorbox}
\section[être / ser · estar]{être / ser · estar — the same two verbs, twice}

\label{lesson:FR-C16-ser-estar}



\begin{quote}
\emph{Stāre} should have rung a bell. If you have ever looked at Spanish, you have met it wearing its own name.
\end{quote}

\begin{culture}
Spanish has \textbf{two} verbs for \textquotedblleft{}to be,\textquotedblright{} and choosing between them is famously the hardest thing about beginner Spanish:

\begin{itemize}
\raggedright
  \item \textbf{ser} $\leftarrow$ \emph{esse}
  \item \textbf{estar} $\leftarrow$ \emph{stāre}
\end{itemize}

Those are the \textbf{same two Latin verbs} French has. French did not get a different inheritance; it made a different decision with the same one.

\begin{itemize}
\raggedright
  \item \textbf{Spanish} kept them \textbf{apart} as two verbs, and made every speaker choose between them for the rest of their life.
  \item \textbf{French} kept \textbf{one} verb, \emph{être} $\leftarrow$ \emph{esse}, and quietly \textbf{swallowed} a large limb of \emph{stāre} — the \emph{ét-} forms — into it.
\end{itemize}

Neither language invented anything. They took the same two words and resolved the overlap in opposite directions, and the result is that a French speaker never thinks about \emph{stāre} while a Spanish speaker thinks about it every day.

This is worth carrying beyond this chapter, because it is the general shape of how related languages differ. They rarely differ because one has something the other lacks. They differ because both inherited the same crowded set of options and \textbf{tidied it up differently}.
\end{culture}

\subsection*{Your turn}
Say these aloud:

\begin{itemize}
\raggedright
  \item the Spanish pair — \textquotedblleft{}ser, estar\textquotedblright{}
  \item their sources — \textquotedblleft{}ser from esse, estar from stāre\textquotedblright{}
  \item the French answer — \textquotedblleft{}one verb, être, which absorbed the ét- limb\textquotedblright{}
  \item the finding — \textquotedblleft{}same two words, opposite solutions\textquotedblright{}
\end{itemize}

\begin{tcolorbox}[breakable,colback=teal!4,colframe=teal!35!black,title={Before you move on}]
Which two Latin verbs lie behind Spanish \emph{ser} and \emph{estar}? (\emph{\textbf{esse}} and \emph{\textbf{stāre}}.) Are they the same two behind French \emph{être}? (\textbf{Yes} — the identical inheritance.) What did each language do? (\textbf{Spanish kept them apart} as two verbs; \textbf{French kept one} and absorbed \emph{stāre}'s \emph{ét-} forms.) Why do related languages usually differ? (They inherited the \textbf{same options} and tidied them differently.) Next: all six forms, together.
\end{tcolorbox}
\section[Practice]{Practice — the whole of être}

\label{lesson:FR-C16-practice}



\begin{quote}
Six forms and three ancestors. Run the forms first; the history is what makes them stay.
\end{quote}

\subsection*{Your turn: the six}
Say these aloud:

\begin{itemize}
\raggedright
  \item \textquotedblleft{}je suis, tu es, il est, nous sommes, vous êtes, ils sont\textquotedblright{}
  \item backwards — \textquotedblleft{}ils sont, vous êtes, nous sommes, il est, tu es, je suis\textquotedblright{}
  \item two sentences — \textquotedblleft{}Je suis français\textquotedblright{}, \textquotedblleft{}Je suis fatigué\textquotedblright{}
\end{itemize}

\emph{Twice through:} \textquotedblleft{}je suis, tu es, il est, nous sommes, vous êtes, ils sont\textquotedblright{}

\subsection*{Your turn: the two verbs side by side}
Say these aloud:

\begin{itemize}
\raggedright
  \item the buzz and the hiss — \textquotedblleft{}ils ont, ils sont\textquotedblright{}
  \item the vowel pair — \textquotedblleft{}tu as, tu es\textquotedblright{}
  \item both verbs from the top — \textquotedblleft{}j'ai … je suis\textquotedblright{}
\end{itemize}

\subsection*{Your turn: the three stems}
Say these aloud:

\begin{itemize}
\raggedright
  \item \textquotedblleft{}es- from esse: the present and the infinitive\textquotedblright{}
  \item \textquotedblleft{}fu- from fuī: the passé simple, and English \emph{be}\textquotedblright{}
  \item \textquotedblleft{}ét- from stāre: to stand\textquotedblright{}
  \item the name — \textquotedblleft{}suppletion\textquotedblright{}; and \textquotedblleft{}go, went\textquotedblright{}
\end{itemize}

\begin{grammarlens}[title={What you've built this chapter}]
\begin{itemize}
\raggedright
  \item \textbf{the second verb the language rests on}, six forms, one lesson each.
  \item \textbf{an explanation instead of a pattern.} There is no pattern in \emph{être} because it is \textbf{three verbs}: \emph{esse}, \emph{fuī} and \emph{stāre}, patched into one paradigm.
  \item \textbf{a word for that} — \textbf{suppletion} — and the knowledge that English does exactly the same with \emph{go}/\emph{went} and \emph{am}/\emph{be}/\emph{was}.
  \item \textbf{a comparison that generalises.} Spanish kept \emph{ser} and \emph{estar} apart; French kept one verb and absorbed the other's limb. Related languages usually differ not by what they inherited but by \textbf{how they tidied it}.
\end{itemize}
\end{grammarlens}

\begin{tcolorbox}[breakable,colback=teal!4,colframe=teal!35!black,title={Before you move on}]
Give all six forms of \emph{être}. How do you hear \emph{ils sont} against \emph{ils ont}? (\textbf{Hiss} against \textbf{buzz}.) Name the three stems and their sources. (\textbf{es-} $\leftarrow$ \emph{esse}, \textbf{fu-} $\leftarrow$ \emph{fuī}, \textbf{ét-} $\leftarrow$ \emph{stāre}.) What is the process called? (\textbf{Suppletion}.) Which Spanish verb is \emph{stāre}? (\emph{\textbf{estar}}.)
\end{tcolorbox}
